% =====================================================================
\chapter{Contrastes de hipótesis}
\label{tema:5}
% =====================================================================

Un contraste de hipótesis decide, a partir de los datos, si hay evidencia suficiente contra una afirmación sobre la población: que una moneda es equilibrada, que un nuevo modelo no mejora al anterior o que unos recuentos siguen una distribución dada. Este tema presenta el marco de Neyman y Pearson (hipótesis nula y alternativa, región de rechazo, errores de tipo~I y~II, potencia), el $p$-valor, los contrastes $z$ y $t$ para medias y proporciones, el contraste $\chi^{2}$ de bondad de ajuste y el problema de las comparaciones múltiples, con su aplicación a la evaluación de modelos de lenguaje en muchos benchmarks. Se sigue el capítulo~10 de Wasserman~\cite{wasserman}, «Hypothesis Testing and p-values», las lecturas~17b, 18 y~19 del curso 18.05 del MIT~\cite{mit1805}, las clases~5, 7 y~11 de Stats~200~\cite{stats200} y los capítulos~9 y~11 de OpenStax~\cite{openstax}.

\section{Planteamiento de un contraste}

\begin{definicion}[Hipótesis, estadístico y región de rechazo]\label{def:t5-contraste}
Sea $\{F_\theta:\theta\in\Theta\}$ un modelo y $\Theta_0\subset\Theta$. Se contrastan la \emph{hipótesis nula} $H_0:\theta\in\Theta_0$ frente a la \emph{hipótesis alternativa} $H_1:\theta\in\Theta\setminus\Theta_0$. Una hipótesis es \emph{simple} si especifica un único valor de $\theta$ y \emph{compuesta} en otro caso. Un \emph{contraste} consiste en un \emph{estadístico de contraste} $T=T(X_1,\dots,X_n)$ y una \emph{región de rechazo} $R$: si $T\in R$ se rechaza $H_0$; si no, no se rechaza.
\end{definicion}

La hipótesis nula representa la situación por defecto («no hay efecto», «la moneda es equilibrada», «los dos modelos son igual de buenos») y el contraste busca evidencia contra ella. Las alternativas más comunes para un parámetro real son la \emph{bilateral} $H_1:\theta\neq\theta_0$ y las \emph{unilaterales} $H_1:\theta>\theta_0$ o $H_1:\theta<\theta_0$.

\begin{definicion}[Errores, nivel y potencia]\label{def:t5-errores}
Se comete un \emph{error de tipo~I} si se rechaza $H_0$ siendo cierta, y un \emph{error de tipo~II} si no se rechaza siendo falsa:
\begin{center}
\small
\begin{tabular}{@{}lcc@{}}
\toprule
 & $H_0$ cierta & $H_0$ falsa \\
\midrule
No se rechaza $H_0$ & decisión correcta & error de tipo~II \\
Se rechaza $H_0$ & error de tipo~I & decisión correcta \\
\bottomrule
\end{tabular}
\end{center}
El contraste tiene \emph{nivel de significación} $\alpha$ si $\PP_\theta(T\in R)\le\alpha$ para todo $\theta\in\Theta_0$. La \emph{función de potencia} es $\pi(\theta)=\PP_\theta(T\in R)$; para $\theta\in\Theta\setminus\Theta_0$, $\pi(\theta)$ es la \emph{potencia} contra esa alternativa y $\beta(\theta)=1-\pi(\theta)$ es la probabilidad de error de tipo~II~\cite[cap.~10]{wasserman}.
\end{definicion}

Los dos errores no se pueden reducir a la vez con un tamaño muestral fijo: agrandar la región de rechazo reduce los errores de tipo~II y aumenta los de tipo~I. El convenio es fijar $\alpha$ (habitualmente $0{,}05$ o $0{,}01$) y, entre los contrastes de nivel $\alpha$, preferir los más potentes. Las dos hipótesis no son simétricas: no rechazar $H_0$ no prueba que sea cierta, solo que los datos no aportan evidencia suficiente contra ella.

\section{El \texorpdfstring{$p$}{p}-valor}

\begin{definicion}[$p$-valor]\label{def:t5-pvalor}
Supóngase que se rechaza $H_0$ para valores grandes de un estadístico $T$ y que $H_0$ es simple. Si el valor observado es $t_{\mathrm{obs}}$, el \emph{$p$-valor} es
\[
p=\PP_{H_0}(T\ge t_{\mathrm{obs}}),
\]
la probabilidad, calculada suponiendo $H_0$ cierta, de obtener un resultado al menos tan extremo como el observado. (Si se rechaza para valores pequeños, $p=\PP_{H_0}(T\le t_{\mathrm{obs}})$; en un contraste bilateral con estadístico de distribución simétrica, $p=\PP_{H_0}(\abs T\ge\abs{t_{\mathrm{obs}}})$.)
\end{definicion}

\begin{proposicion}[Decisión con el $p$-valor]\label{prop:t5-pvalor}
Sea un contraste que rechaza $H_0$ si $T\ge c_\alpha$, donde $c_\alpha$ cumple $\PP_{H_0}(T\ge c_\alpha)=\alpha$ y $T$ tiene distribución continua bajo $H_0$. Entonces se rechaza $H_0$ al nivel $\alpha$ si y solo si $p\le\alpha$.
\end{proposicion}

El $p$-valor es, por tanto, el menor nivel al que los datos permiten rechazar $H_0$, y resume en un número la fuerza de la evidencia contra $H_0$: cuanto menor, más incompatibles son los datos con $H_0$.

\begin{observacion}[Lo que el $p$-valor no es]\label{obs:t5-pvalor}
El $p$-valor \emph{no} es la probabilidad de que $H_0$ sea cierta, ni la probabilidad de que los resultados se deban al azar, ni mide el tamaño del efecto. Un $p$-valor muy pequeño puede acompañar a un efecto irrelevante si $n$ es enorme, y un efecto grande puede dar un $p$-valor alto si $n$ es pequeño. Por eso conviene informar siempre del efecto estimado con su intervalo de confianza, además del $p$-valor~\cite[lecturas~17b y~20]{mit1805}.
\end{observacion}

\section{Contrastes \texorpdfstring{$z$}{z} y \texorpdfstring{$t$}{t}}

Los contrastes de esta sección usan como estadístico una estimación tipificada con su error estándar calculado bajo $H_0$. En todos ellos, las regiones de rechazo de nivel $\alpha$ son $T>q_{1-\alpha}$ (alternativa $>$), $T<-q_{1-\alpha}$ (alternativa $<$) y $\abs T>q_{1-\alpha/2}$ (bilateral), donde $q$ es el cuantil de la distribución de referencia.

\begin{center}
\small
\begin{tabular}{@{}llll@{}}
\toprule
Contraste & $H_0$ & Estadístico & Distribución bajo $H_0$ \\
\midrule
$z$ para la media ($\sigma$ conocida) & $\mu=\mu_0$ & $Z=\dfrac{\Xbar-\mu_0}{\sigma/\sqrt n}$ & $\Nor(0,1)$ (datos normales) \\[2ex]
$t$ para la media & $\mu=\mu_0$ & $T=\dfrac{\Xbar-\mu_0}{S/\sqrt n}$ & $t_{n-1}$ (datos normales) \\[2ex]
$z$ para una proporción & $p=p_0$ & $Z=\dfrac{\widehat p-p_0}{\sqrt{p_0(1-p_0)/n}}$ & $\approx\Nor(0,1)$ ($n$ grande) \\
\bottomrule
\end{tabular}
\end{center}

La distribución bajo $H_0$ se sigue del \cref{teo:t3-suma-normales}, del \cref{cor:t3-t} y del \cref{cor:t3-binomial}, respectivamente. En el contraste de una proporción, el error estándar se calcula con $p_0$, no con $\widehat p$, porque se razona bajo $H_0$.

\addtocounter{definicion}{1}%

\section{Potencia y tamaño muestral}

\begin{proposicion}[Potencia del contraste $z$ unilateral]\label{prop:t5-potencia}
Para el contraste $z$ de $H_0:\mu=\mu_0$ frente a $H_1:\mu>\mu_0$ al nivel $\alpha$ con datos $\Nor(\mu,\sigma^{2})$, la función de potencia es
\[
\pi(\mu)=1-\Phi\Bigl(z_{1-\alpha}-\frac{(\mu-\mu_0)\sqrt n}{\sigma}\Bigr).
\]
\end{proposicion}

En el \cref{ej:t5-z}, si la media real fuera $220$, la potencia sería $1-\Phi(1{,}645-2)=\Phi(0{,}355)\approx0{,}64$: con $64$ respuestas, el contraste detectaría esa diferencia solo en dos de cada tres ocasiones. El \cref{eje:t5-potencia} deduce el tamaño muestral necesario para alcanzar una potencia dada.

\section{Contraste \texorpdfstring{$\chi^{2}$}{chi cuadrado} de bondad de ajuste}

Sea una variable categórica con $k$ categorías y probabilidades $p_1,\dots,p_k$. En $n$ observaciones independientes se obtienen los recuentos $O_1,\dots,O_k$, con $\sum_jO_j=n$. Se quiere contrastar $H_0:p_j=p_{0j}$ para todo $j$, con $p_{0j}$ dados.

\begin{definicion}[Estadístico $\chi^{2}$ de Pearson]\label{def:t5-chi2}
Sean $E_j=np_{0j}$ los recuentos \emph{esperados} bajo $H_0$. El estadístico es
\[
X^{2}=\sum_{j=1}^{k}\frac{(O_j-E_j)^{2}}{E_j}.
\]
Bajo $H_0$, y si todos los $E_j$ son al menos~$5$, $X^{2}$ sigue aproximadamente una distribución $\chi^{2}_{k-1}$. Se rechaza $H_0$ al nivel $\alpha$ si $X^{2}>\chi^{2}_{k-1;1-\alpha}$~\cite[§11.2]{openstax}.
\end{definicion}

Los $k-1$ grados de libertad se deben a la restricción $\sum_jO_j=n$: conocidos $k-1$ recuentos, el último queda determinado. Si se estiman $m$ parámetros del modelo a partir de los mismos datos, se restan también: $k-1-m$ grados de libertad. Si algún $E_j$ es menor que $5$, se agrupan categorías contiguas.

\addtocounter{definicion}{1}%

\section{Comparaciones múltiples}

Si se realizan muchos contrastes, cada uno al nivel $\alpha$, la probabilidad de cometer \emph{algún} error de tipo~I crece rápidamente.

\begin{proposicion}\label{prop:t5-multiples}
Si se realizan $m$ contrastes independientes al nivel $\alpha$ y todas las hipótesis nulas son ciertas, la probabilidad de rechazar al menos una es $1-(1-\alpha)^{m}$, y el número esperado de rechazos (todos falsos) es $m\alpha$.
\end{proposicion}

Con $m=20$ y $\alpha=0{,}05$, la probabilidad de al menos un «descubrimiento» falso es $1-0{,}95^{20}\approx0{,}64$.

\begin{definicion}[Tasa de error por familia]\label{def:t5-fwer}
La \emph{tasa de error por familia} (FWER, \emph{family-wise error rate}) de un procedimiento que realiza $m$ contrastes es la probabilidad de cometer al menos un error de tipo~I~\cite[clase~11]{stats200}.
\end{definicion}

\begin{teorema}[Corrección de Bonferroni]\label{teo:t5-bonferroni}
Si cada uno de los $m$ contrastes se realiza al nivel $\alpha/m$ (es decir, se rechaza la hipótesis $i$ si $p_i\le\alpha/m$), la tasa de error por familia es a lo sumo $\alpha$, cualquiera que sea la dependencia entre los contrastes y cualquiera que sea el conjunto de hipótesis nulas ciertas.
\end{teorema}

La corrección de Bonferroni es sencilla y muy general, pero conservadora: con $m$ grande reduce mucho la potencia. Existen procedimientos más potentes, como el de Holm o el de Benjamini y Hochberg, que controla una cantidad distinta, la tasa de falsos descubrimientos~\cite[clase~11]{stats200}; quedan fuera de este módulo.

\begin{observacion}[Comparar modelos en muchos benchmarks]\label{obs:t5-benchmarks}
La evaluación de modelos de lenguaje es un terreno abonado para las comparaciones múltiples: se comparan dos modelos en decenas de benchmarks, se prueban muchas variantes de \emph{prompt} y se informa de la mejor, o se ordena una tabla con cientos de modelos. Si en $20$ benchmarks independientes no hay ninguna diferencia real, se esperan $20\cdot0{,}05=1$ diferencias «significativas» al $5\,\%$. Para comparar la exactitud de dos modelos se puede usar un contraste $z$ de diferencia de proporciones (\cref{eje:t5-dos-modelos}); cuando ambos modelos responden a las \emph{mismas} preguntas, es preferible un contraste pareado que use las diferencias pregunta a pregunta, porque elimina la variabilidad debida a la dificultad de cada pregunta (ampliación).
\end{observacion}

\addtocounter{definicion}{1}%

\addtocounter{definicion}{1}%

